% Modèle matOS : sujet d'examen avec barème. Compiler DEUX fois avec pdflatex
% (le total des points est calculé à la seconde compilation).
% Corrigé : remplacer [addpoints] par [addpoints,answers].
\documentclass[11pt,a4paper,addpoints]{exam}
\usepackage[T1]{fontenc}
\usepackage[utf8]{inputenc}
\usepackage[french]{babel}
\usepackage{lmodern,microtype}
\usepackage{amsmath,amssymb,mathtools}

\pointname{ pt}
\pointsinrightmargin
\renewcommand{\solutiontitle}{\noindent\textbf{Corrigé.}\enspace}
\newcommand{\R}{\mathbb{R}}

\pagestyle{headandfoot}
\firstpageheader{Université d'Antananarivo\\Mention Mathématiques}{}{L1 --- Analyse 1\\Session de janvier 2027}
\runningheader{Analyse 1}{}{Session de janvier 2027}
\footer{}{Page \thepage\ sur \numpages}{}

\begin{document}
\begin{center}
  {\Large\bfseries Examen d'analyse 1}\\[4pt]
  Durée : 2 heures. Documents et calculatrices interdits.\\
  Le sujet comporte \numquestions\ exercices, notés sur \numpoints\ points.
\end{center}
\bigskip

\begin{questions}

\question[6] Questions de cours.
\begin{parts}
  \part[2] Énoncer le théorème des valeurs intermédiaires.
  \part[4] Démontrer qu'une suite convergente est bornée.
  \begin{solution}
  Si $u_n \to \ell$, il existe $N$ tel que $|u_n - \ell| \leq 1$ pour $n \geq N$ ;
  alors $|u_n| \leq \max(|u_0|, \dots, |u_{N-1}|, |\ell| + 1)$.
  \end{solution}
\end{parts}

\question[8] Soit $f : \R \to \R$ définie par $f(x) = x^3 - 3x + 1$.
\begin{parts}
  \part[3] Dresser le tableau de variations de $f$.
  \part[5] Montrer que l'équation $f(x) = 0$ admet exactement trois solutions réelles.
\end{parts}

\question[6] Calculer $\displaystyle \int_0^1 x\,e^{x}\,dx$.
\begin{solution}
Par parties : $\bigl[x e^x\bigr]_0^1 - \int_0^1 e^x\,dx = e - (e - 1) = 1$.
\end{solution}

\end{questions}
\end{document}
